 \section{聖經中耶穌的講論（按照時間順序分）}
 
 \subsection{初期的講論}
 \begin{table}[htbp]
  \centering

    \begin{tabular}{p{2.43em}|p{3.93em}|p{11.53em}|c|p{4.045em}|p{4.045em}|p{4.045em}|p{4.045em}}
    \toprule
    人群    & \multicolumn{2}{p{16.5em}|}{事 件} & 地 點 & 馬太福音 & 馬可福音 & 路加福音 & 約翰福音 \\    \midrule

\multirow{2}{*}{個人} & \multicolumn{2}{p{13.5em}|}{與尼哥底母談道（逾越節）} & 耶路撒冷 &       &       &       & 3:1-21 \\
    
\cmidrule{2-8}       & \multicolumn{2}{p{11.5em}|}{與撒馬利亞婦人傳道} & 敘加 &       &       &       & 4:3下-38 \\ \midrule


\multirow{1}{*}{眾人}&\multicolumn{2}{p{16.5em}|}{約翰下監安息日讲撒勒法寡婦和乃縵} & 拿撒勒會堂 &4:12-13a&&4:16-30&\\\hline

\cmidrule{1-8}   \multirow{6}[10]{*}{眾人} & \multicolumn{2}{p{13.15em}|}{耶穌在馬太家坐席論門徒禁食} &       & 9:10-17 & 2:15-22 & 5:29-39 &  \\

\cmidrule{2-8}          & \multirow{4}{*}{\parbox{0.08\textwidth}{耶穌講論父與子的關係（逾越節）}} & a.子照著父所作的事去行 & \multirow{4}{*}{耶路撒冷} &       &       &       & 5:17-21 \\

\cmidrule{3-3}\cmidrule{5-8}      &       & b.父把審判的事交給子 &       &       &       &       & 5:22-30 \\
\cmidrule{3-3}\cmidrule{5-8}          &       & c.父爲子作見證 &       &       &       &       & 5:31-42 \\
\cmidrule{3-3}\cmidrule{5-8}          &       & d.子奉父的命而來 &       &       &       &       & 5:43-47 \\

\cmidrule{1-8} &\multicolumn{2}{p{12.em}|}{門徒安息日掐麥穗引辯論}&\multirow{1}{*}{從麥地經過} &12:1-8&2:23-28&6:1-5& \\ \bottomrule
    \end{tabular}%
  \label{tab:addlabel}%
    \caption{初期的講論}
\end{table}%
 

 
 
%\subsection{到第三個逾越節之前的講論} 
 \subsection{在迦百農設立十二門徒後的講論}
\begin{table}[H]
\centering 
\begin{tabular}{p{13.07em}|p{8.43em}|p{5.285em}|p{5.57em}|p{5.285em}|p{4.855em}}
\toprule
\textbf{事 件} & \textbf{地 點} & \textbf{太} & \textbf{可} & \textbf{路} & \textbf{約} \\\midrule


登山教訓門徒論八福& 近迦百農山上  & 5:1-12& \multicolumn{1}{r|}{} & \multicolumn{1}{l|}{6:20-23} & \multicolumn{1}{r}{} \\    \midrule

富足飽足喜笑好人的四禍& 近迦百農山上&  & \multicolumn{1}{r|}{} & \multicolumn{1}{l|}{6:24-26} & \multicolumn{1}{r}{} \\    \midrule

門徒為鹽為光& 近迦百農山上& 5:13-16& \multicolumn{1}{r|}{} & \multicolumn{1}{r|}{} & \multicolumn{1}{r}{} \\    \midrule

成全律法和先知的道&\multicolumn{1}{l|}{\multirow{6}{*}{律法和先知}}& 5:17-20& \multicolumn{1}{r|}{} & \multicolumn{1}{r|}{} & \multicolumn{1}{r}{} \\   
\cmidrule{1-1}\cmidrule{3-6}{不可殺人愛仇敵}&& 5:21-26& \multicolumn{1}{r|}{} & \multicolumn{1}{r|}{} & \multicolumn{1}{r}{} \\   
\cmidrule{1-1}\cmidrule{3-6}{不可姦淫休妻}&   & 5:27-32& \multicolumn{1}{r|}{} & \multicolumn{1}{r|}{} & \multicolumn{1}{r}{} \\    
\cmidrule{1-1}\cmidrule{3-6}{不可背誓}&   & 5:33-37& \multicolumn{1}{r|}{} & \multicolumn{1}{r|}{} & \multicolumn{1}{r}{} \\   
\cmidrule{1-1}\cmidrule{3-6}{以眼還眼，以牙還牙}&   & 5:38-42& \multicolumn{1}{r|}{} & \multicolumn{1}{r|}{} & \multicolumn{1}{r}{} \\   
\cmidrule{1-1}\cmidrule{3-6}{愛鄰舍恨仇敵}&  & 5:43-48& \multicolumn{1}{r|}{} & \multicolumn{1}{l|}{6:27-38} & \multicolumn{1}{r}{} \\    \midrule

施捨,禱告,禁食&\multicolumn{1}{l|}{\multirow{1}{*}{行在暗中}}& 6:1-18& \multicolumn{1}{r|}{} & \multicolumn{1}{r|}{} & \multicolumn{1}{r}{} \\  \midrule


積攢財寶在天上& 近迦百農  & 6:19-21& \multicolumn{1}{r|}{} & \multicolumn{1}{r|}{} & \multicolumn{1}{r}{} \\    \midrule
眼睛就是身上的燈& 近迦百農  & 6:22-23 & \multicolumn{1}{r|}{} & \multicolumn{1}{r|}{} & \multicolumn{1}{r}{} \\    \midrule
不能侍奉兩個主& 近迦百農  & 6:24 & \multicolumn{1}{r|}{} & \multicolumn{1}{r|}{} & \multicolumn{1}{r}{} \\    \midrule
勿慮衣食不為明天憂慮& 近迦百農  & 6:25-34 & \multicolumn{1}{r|}{} & \multicolumn{1}{r|}{} & \multicolumn{1}{r}{} \\    \midrule

去掉自己眼中的梁木& 近迦百農  &  & \multicolumn{1}{r|}{} & \multicolumn{1}{l|}{6:39-42} & \multicolumn{1}{r}{} \\    \midrule

毋論斷人& 近迦百農  & 7:1-5 & \multicolumn{1}{r|}{} & \multicolumn{1}{r|}{} & \multicolumn{1}{r}{} \\    \midrule
不要把聖物給狗豬& 近迦百農  & 7:6 & \multicolumn{1}{r|}{} & \multicolumn{1}{r|}{} & \multicolumn{1}{r}{} \\    \midrule
祈求尋找叩門& 近迦百農  & 7:7-12 & \multicolumn{1}{r|}{} & \multicolumn{1}{r|}{} & \multicolumn{1}{r}{} \\    \midrule
兩條門路& 近迦百農  & 7:13-14 & \multicolumn{1}{r|}{} & \multicolumn{1}{r|}{} & \multicolumn{1}{r}{} \\    \midrule
兩種果樹,兩等根基& 近迦百農  & 7:15-27 & \multicolumn{1}{r|}{} & \multicolumn{1}{l|}{6:43-49} & \multicolumn{1}{r}{} \\    \midrule

\textcolor{blue}{差人走后主對眾人讚施洗約翰}& 加利利   & 11:2-15& \multicolumn{1}{r|}{} & 7:18-29& \multicolumn{1}{r}{} \\ \midrule

警法利賽人律法師辨别世代& 加利利   & 11:16-19& \multicolumn{1}{r|}{} & 7:30-35& \multicolumn{1}{r}{} \\ \midrule

%耶穌坐席时有罪女人膏主 & 迦百農   & \multicolumn{1}{r|}{} & \multicolumn{1}{r|}{} & 7:36-50 & \multicolumn{1}{r}{} \\\midrule
\end{tabular}%
\end{table}%


\subsubsection{拣选十二使徒后的初期培训}
\begin{table}[H]
\centering 
\begin{tabular}{p{13.07em}|p{8.43em}|p{5.285em}|p{5.57em}|p{5.285em}|p{4.855em}}
\toprule
\textbf{事 件} & \textbf{地 點} & \textbf{太} & \textbf{可} & \textbf{路} & \textbf{約} \\\midrule

坐船上對眾人講天國結實百倍& 加利利迦百農海邊& \multicolumn{1}{l|}{13:1-9} & \multicolumn{1}{l|}{4:1-9} & 8:4-8 & \multicolumn{1}{r}{} \\ \midrule
對門徒解释用比喻的原因& 加利利迦百農海邊& \multicolumn{1}{l|}{13:10-16} & \multicolumn{1}{l|}{4:10-12} & 8:9-10 & \multicolumn{1}{r}{} \\ \midrule
對門徒解明撒種的比喻& 加利利迦百農海邊& \multicolumn{1}{l|}{13:17-23} & \multicolumn{1}{l|}{4:13-20} & 8:11-15 & \multicolumn{1}{r}{} \\ \midrule

置燈于臺上喻掩藏的事必顯明& 加利利迦百農海邊& \multicolumn{1}{l|}{} & \multicolumn{1}{l|}{4:21-23} & 8:16-17 & \multicolumn{1}{r}{} \\ \midrule

用量器量給人必得回报且多得& 加利利迦百農海邊& \multicolumn{1}{l|}{} & \multicolumn{1}{l|}{4:24} & & \multicolumn{1}{r}{} \\ \midrule

凡有的還要加給他& 加利利迦百農海邊& \multicolumn{1}{l|}{} & \multicolumn{1}{l|}{4:25} & 8:18& \multicolumn{1}{r}{} \\ \midrule

收割時候分开田裡麥子稗子& 加利利迦百農海邊& \multicolumn{1}{l|}{13:24-30} & \multicolumn{1}{l|}{4:26-29} & & \multicolumn{1}{r}{} \\ \midrule

芥菜種宿天上飛鳥的比喻& 加利利迦百農海邊& \multicolumn{1}{l|}{13:31-32} & \multicolumn{1}{l|}{4:30-32} & & \multicolumn{1}{r}{} \\ \midrule

對眾人用比喻给門徒解明比喻& 加利利迦百農海邊& \multicolumn{1}{l|}{} & \multicolumn{1}{l|}{4:33-34} & & \multicolumn{1}{r}{} \\ \midrule

婦人麵酵的比喻& 加利利迦百農海邊& \multicolumn{1}{l|}{13:33} & \multicolumn{1}{l|}{} & & \multicolumn{1}{r}{} \\ \midrule
用比喻對人說話應驗先知预言& 加利利迦百農海邊& \multicolumn{1}{l|}{13:34-35} & \multicolumn{1}{l|}{} & & \multicolumn{1}{r}{} \\ \midrule
對門徒解明稗子的比喻& 加利利迦百農房子& \multicolumn{1}{l|}{13:36-43} & \multicolumn{1}{l|}{} & & \multicolumn{1}{r}{} \\ \midrule
變賣一切買藏寶之地& 加利利迦百農海邊& \multicolumn{1}{l|}{13:44}& \multicolumn{1}{l|}{} && \multicolumn{1}{r}{} \\ \midrule
買賣人變賣一切買好珠子& 加利利迦百農海邊& \multicolumn{1}{l|}{13:45-46} & \multicolumn{1}{l|}{} &  & \multicolumn{1}{r}{} \\ \midrule

以天國網撒喻世界末了& 加利利迦百農海邊& \multicolumn{1}{l|}{13:47-50} & \multicolumn{1}{l|}{} && \multicolumn{1}{r}{} \\ \midrule
家主庫裡拿出的新舊東西& 加利利迦百農海邊& \multicolumn{1}{l|}{13:51-52} & \multicolumn{1}{l|}{} & & \multicolumn{1}{r}{} \\ \midrule

耶穌論信者皆为親屬 & 加利利   & 12:46-50 & 3:31-35 & 8:19-21 & \multicolumn{1}{r}{} \\ \midrule

\end{tabular}%
\end{table}%

\subsubsection{差遣十二使徒传道的培训}
\begin{table}[H]
\centering 
\begin{tabular}{p{17.07em}|p{8.43em}|p{5.285em}|p{5em}|p{4.285em}|p{1.855em}}\toprule
\textbf{事 件} & \textbf{地 點} & \textbf{太} & \textbf{可} & \textbf{路} & \textbf{約} \\\midrule

%\textcolor{blue}{給十二使徒權柄兩兩出去傳道} & 加利利   & 10:1& 6:7& 9:1-2& \multicolumn{1}{r}{} \\ \midrule
%\textbf{往以色列傳天國近了} &地点/内容& 10:5-6 &  &  & \multicolumn{1}{r}{} \\ \midrule

\textbf{免费醫病潔淨痲瘋趕鬼} &方法& 10:8& & & \multicolumn{1}{r}{} \\ \midrule
\textbf{不帶食物錢靠做工得食} & 原则& 10:9-10& 6:8-9& 9:3 & \multicolumn{1}{r}{} \\ \midrule
\textbf{進谁家就住在那裡/如同羊進入狼群} &住宿/危险 & 10:11-15& 6:10-11& 9:4-5 & \multicolumn{1}{r}{} \\ \midrule

\textbf{忍耐到底必然得救/不要懼怕} &态度& 10:21-23&& & \multicolumn{1}{r}{} \\ \midrule

\textbf{背著他的十字架/接待義人得義的賞賜} &代价/賞賜&10:34-39&  && \multicolumn{1}{r}{} \\ \midrule
\end{tabular}%
\end{table}%

\subsubsection{施洗约翰被杀后的工作}
\begin{table}[H]\centering 
\begin{tabular}{p{13.07em}|p{8.43em}|p{5.285em}|p{5.57em}|p{5.285em}|p{4.855em}}
\toprule
\textbf{事 件} & \textbf{地 點} & \textbf{太} & \textbf{可} & \textbf{路} & \textbf{約} \\\midrule

論信神所差來的就是做神的工& 迦百農海邊& \multicolumn{1}{r|}{} & \multicolumn{1}{r|}{} &       & \multicolumn{1}{p{4.855em}}{6:22-29} \\\midrule
耶穌論生命的糧 & 迦百農會堂   & \multicolumn{1}{r|}{} & \multicolumn{1}{r|}{} &       & \multicolumn{1}{p{4.855em}}{6:30-40} \\\midrule
對猶太人論吃人子肉喝人子血& 迦百農會堂   & \multicolumn{1}{r|}{} & \multicolumn{1}{r|}{} &       & \multicolumn{1}{p{4.855em}}{6:41-59} \\\midrule

對門徒論話就是靈就是生命& 迦百農& \multicolumn{1}{r|}{} & \multicolumn{1}{r|}{} &       & \multicolumn{1}{p{4.855em}}{6:60-66} \\\midrule

俗手吃饭和出口的汙穢人&  & 15:1-11& 7:1-16&       &  \\\midrule
%主责应验以赛亚"嘴唇敬主"&  & 15:7-9& 7:6-7 &       &  \\\midrule
%主责其廢神誡命守人遺傳&  & 15:3-6& 7:8-13 &       &  \\\midrule
%對眾人說出口的乃能汙穢人&  & 15:10-11 & 7:14-16&       &  \\\midrule

%對彼得解释出口的汙穢人& 離開眾人進了屋子& 15:15-20& 7:17-23&       &  \\\midrule


責法利賽撒都該人求天上神迹& \multicolumn{1}{p{7.93em}|}{馬加丹/大瑪努他} & \multicolumn{1}{p{5.285em}|}{16:1-4} & \multicolumn{1}{p{5.57em}|}{8:11-13} & \multicolumn{1}{r|}{} &  \\ \midrule

防法利賽/撒都該/希律的酵& \multicolumn{1}{p{7.93em}|}{加利利海船上} & \multicolumn{1}{p{5.285em}|}{16:5-12} & \multicolumn{1}{p{5.57em}|}{8:14-21} & \multicolumn{1}{r|}{} &  \\ \midrule

\end{tabular}%
%\caption{到第三個逾越節}
\end{table}%



\subsection{启示自己受死并复活}
\begin{table}[H] \centering 
\begin{tabular}{p{13.07em}|l|l|l|p{5.285em}|r}\toprule
\textbf{事 件} & \multicolumn{1}{p{7.93em}|}{\textbf{地 點}} & \multicolumn{1}{p{5.285em}|}{\textbf{馬太福音}} & \multicolumn{1}{p{5.57em}|}{\textbf{馬可福音}} & \textbf{路加福音} & \multicolumn{1}{p{4.855em}}{\textbf{約翰福音}} \\\midrule

%途中主問門徒:人說人子是誰& \multicolumn{1}{p{7.93em}|}{該撒利亞腓立比} & \multicolumn{1}{p{5.285em}|}{16:13-14} & \multicolumn{1}{p{5.57em}|}{8:27-28} & 9:18-19&  \\ \midrule

%彼得認耶穌爲神的兒子 & \multicolumn{1}{p{7.93em}|}{該撒利亞腓立比} & \multicolumn{1}{p{5.285em}|}{16:15-16} & \multicolumn{1}{p{5.57em}|}{8:29-30} & 9:20-21&  \\ \midrule

对彼得论把教會建造在磐石上& \multicolumn{1}{p{7.93em}|}{該撒利亞腓立比} & \multicolumn{1}{p{5.285em}|}{16:17-18a} & \multicolumn{1}{p{5.57em}|}{} &  &  \\ \midrule

天國鑰匙/捆綁釋放權柄& \multicolumn{1}{p{7.93em}|}{該撒利亞腓立比} & \multicolumn{1}{p{5.285em}|}{16:18b-20} & \multicolumn{1}{p{5.57em}|}{} &  &  \\ \midrule

\textcolor{red}{首次}提說在耶路撒冷被殺復活& \multicolumn{1}{p{7.93em}|}{該撒利亞腓立比} & \multicolumn{1}{p{5.285em}|}{16:21} & \multicolumn{1}{p{5.57em}|}{8:31} & 9:22&  \\ \midrule



教導門徒背十字架跟從主& \multicolumn{1}{p{7.93em}|}{該撒利亞腓立比} & \multicolumn{1}{p{5.285em}|}{16:24-26} & \multicolumn{1}{p{5.57em}|}{8:34-37} & 9:23-25&  \\ \midrule


人子在榮耀降臨時不认背主者& \multicolumn{1}{p{7.93em}|}{該撒利亞腓立比} & \multicolumn{1}{p{5.285em}|}{16:27} & \multicolumn{1}{p{5.57em}|}{8:38} & 9:26&  \\ \midrule



以約翰譬以利亞/\textcolor{red}{人子要受害}& \multicolumn{1}{p{7.93em}|}{下腓立比黑門山} & \multicolumn{1}{p{5.285em}|}{17:10-13} & \multicolumn{1}{p{5.57em}|}{9:11-13} &  &  \\ \midrule




\textbf{耶穌對門徒\textcolor{red}{三論}受害的事} & \multicolumn{1}{p{7.93em}|}{住在加利利} & \multicolumn{1}{p{5.285em}|}{17:22-23} & \multicolumn{1}{p{5.57em}|}{9:30-32} & 9:43b-45 &  \\ \midrule



願意做首的必做眾人的用人& \multicolumn{1}{p{7.93em}|}{迦百農} & \multicolumn{1}{p{5.285em}|}{} & \multicolumn{1}{p{5.57em}|}{9:35} &  &  \\ \midrule


不成小孩子不得進天國& \multicolumn{1}{p{7.93em}|}{迦百農在屋裡} & \multicolumn{1}{p{5.285em}|}{18:2-3} & \multicolumn{1}{p{5.57em}|}{10:15} & &  \\ \midrule

謙卑像這小孩子是最大的& \multicolumn{1}{p{7.93em}|}{迦百農在屋裡} & \multicolumn{1}{p{5.285em}|}{18:4} & \multicolumn{1}{p{5.57em}|}{} & &  \\ \midrule


接待小孩就是接待主& \multicolumn{1}{p{7.93em}|}{迦百農} & \multicolumn{1}{p{5.285em}|}{18:5} & \multicolumn{1}{p{5.57em}|}{9:36-37} & 9:47-48a &  \\ \midrule

你們中間最小的他便為大& \multicolumn{1}{p{7.93em}|}{迦百農} & \multicolumn{1}{p{5.285em}|}{} & \multicolumn{1}{p{5.57em}|}{9:36-37} & 9:48b &  \\ \midrule



教導約翰不要禁止人趕鬼& \multicolumn{1}{p{7.93em}|}{迦百農} &&\multicolumn{1}{p{5.57em}|}{9:38-40} & 9:49-50 &  \\\midrule

因你們屬基督給水喝必得賞賜& \multicolumn{1}{p{7.93em}|}{迦百農} &&\multicolumn{1}{p{5.57em}|}{9:41} &  &  \\\midrule


残肢入永生强如全人进地狱& \multicolumn{1}{p{7.93em}|}{迦百農?} &\multicolumn{1}{p{5.285em}|}{18:8-9}&9:43-47&  &  \\\midrule
地狱蟲不死火不滅鹽醃各人& \multicolumn{1}{p{7.93em}|}{迦百農?} &\multicolumn{1}{p{5.285em}|}{}&9:48-49&  &  \\\midrule
裡頭應當有鹽彼此和睦& \multicolumn{1}{p{7.93em}|}{迦百農?} &\multicolumn{1}{p{5.285em}|}{}&9:50&  &  \\\midrule

小子的使者在天上常見天父& \multicolumn{1}{p{7.93em}|}{迦百農?} &\multicolumn{1}{p{5.285em}|}{18:10}&9:42-50&  &  \\\midrule


被弟兄得罪时的处理方法& \multicolumn{1}{p{7.93em}|}{迦百農?} &\multicolumn{1}{p{5.285em}|}{18:15-17}&&  &  \\\midrule

在地上捆綁和釋放的权柄& \multicolumn{1}{p{7.93em}|}{迦百農?} &\multicolumn{1}{p{5.285em}|}{18:18}&&  &  \\\midrule
两人同心求就有主同在& \multicolumn{1}{p{7.93em}|}{迦百農?} &\multicolumn{1}{p{5.285em}|}{18:19-20}&&  &  \\\midrule
對彼得说饒恕七十個七次& \multicolumn{1}{p{7.93em}|}{迦百農?} &\multicolumn{1}{p{5.285em}|}{18:21-22}&&  &  \\\midrule
欠十兩銀子的仆人& \multicolumn{1}{p{7.93em}|}{迦百農?} &\multicolumn{1}{p{5.285em}|}{18:23-35}&&  &  \\\midrule
\end{tabular}
\end{table}%label{tab:addlabel}%


\subsubsection{耶穌在住棚節時的講論}
\begin{table}[H]\centering 
\begin{tabular}{p{13.07em}|l|l|l|p{5.285em}|r}\toprule
\textbf{事 件} & \multicolumn{1}{p{7.93em}|}{\textbf{地 點}} & \multicolumn{1}{p{5.285em}|}{\textbf{馬太福音}} & \multicolumn{1}{p{5.57em}|}{\textbf{馬可福音}} & \textbf{路加福音} & \multicolumn{1}{p{4.855em}}{\textbf{約翰福音}} \\\midrule 


對跟随文士論人子无枕之地& \multicolumn{1}{p{7.93em}|}{迦百農渡口} & \multicolumn{1}{p{5.285em}|}{8:18-20} &       & 9:57-58&  \\ \midrule

對葬父之人論任死人葬其死人& \multicolumn{1}{p{7.93em}|}{迦百農渡口} & \multicolumn{1}{p{5.285em}|}{8:21-22} &       & 9:59-60&  \\ \midrule

對人論扶犁向後看不配進神國& \multicolumn{1}{p{7.93em}|}{迦百農渡口} & \multicolumn{1}{p{5.285em}|}{} &       & 9:61-62 &  \\ \midrule


\textbf{耶穌求那差他來者的榮耀} & \multicolumn{1}{p{7.93em}|}{住棚節耶路撒冷} &       &       & \multicolumn{1}{r|}{} & \multicolumn{1}{p{4.855em}}{7:16-18} \\\midrule

\textbf{向眾人論安息日治病} & \multicolumn{1}{p{7.93em}|}{住棚節耶路撒冷} &       &       & \multicolumn{1}{r|}{} & \multicolumn{1}{p{4.855em}}{7:19-24} \\\midrule
\textbf{向眾人論父差子} & \multicolumn{1}{p{7.93em}|}{住棚節耶路撒冷} &       &       & \multicolumn{1}{r|}{} & \multicolumn{1}{p{4.855em}}{7:25-29} \\\midrule
\textbf{向眾人論腹中的活水江河} & \multicolumn{1}{p{7.93em}|}{住棚節耶路撒冷} &       &       & \multicolumn{1}{r|}{} & \multicolumn{1}{p{4.855em}}{7:37-44} \\\midrule
    %耶穌赦免犯罪的婦人 & \multicolumn{1}{p{7.93em}|}{住棚節耶路撒冷} &       &       & \multicolumn{1}{r|}{} & \multicolumn{1}{p{4.855em}}{8:1-11} \\ \midrule
    耶穌論光與自由 & \multicolumn{1}{p{7.93em}|}{住棚節耶路撒冷} &       &       & \multicolumn{1}{r|}{} & \multicolumn{1}{p{4.855em}}{8:12-59} \\\midrule
    耶穌醫生來瞎眼者 & \multicolumn{1}{p{7.93em}|}{住棚節耶路撒冷} &       &       & \multicolumn{1}{r|}{} & \multicolumn{1}{p{4.855em}}{9:1-41} \\ \midrule
    耶穌論羊的門好牧人 & \multicolumn{1}{p{7.93em}|}{住棚節耶路撒冷} &       &       & \multicolumn{1}{r|}{} & \multicolumn{1}{p{4.855em}}{10:1-21} \\\midrule 
\end{tabular}%
\end{table}%

\subsubsection{差遣七十人兩兩传道}

\begin{table}[H] \centering 
\begin{tabular}{p{13.07em}|l|l|l|p{5.285em}|r} \toprule
\textbf{事 件} & \multicolumn{1}{p{7.93em}|}{\textbf{地 點}} & \multicolumn{1}{p{5.285em}|}{\textbf{馬太福音}} & \multicolumn{1}{p{5.57em}|}{\textbf{馬可福音}} & \textbf{路加福音} & \multicolumn{1}{p{4.855em}}{\textbf{約翰福音}} \\\midrule 



當求莊稼主打發工人收莊稼& \multicolumn{1}{p{7.93em}|}{} & \multicolumn{1}{p{5.285em}|}{} &       & 10:2&  \\\midrule

如同羊羔進入狼群& \multicolumn{1}{p{7.93em}|}{去耶京途中} & \multicolumn{1}{p{5.285em}|}{} &       & 10:3&  \\\midrule
不帶錢囊口袋鞋& \multicolumn{1}{p{7.93em}|}{去耶京途中？} & \multicolumn{1}{p{5.285em}|}{} &       & 10:4&  \\\midrule
为所住的家求的平安& \multicolumn{1}{p{7.93em}|}{} & \multicolumn{1}{p{5.285em}|}{} &       & 10:5-6&  \\\midrule

靠传福音糊口& \multicolumn{1}{p{7.93em}|}{} & \multicolumn{1}{p{5.285em}|}{} &       & 10:7-8&  \\\midrule
内容：神的國臨近你們了& \multicolumn{1}{p{7.93em}|}{} & \multicolumn{1}{p{5.285em}|}{} &       & 10:9-11&  \\\midrule
哥拉汛伯賽大迦百農之禍& \multicolumn{1}{p{7.93em}|}{} & \multicolumn{1}{p{5.285em}|}{11:20-24} &       & 10:12-15&  \\\midrule
對門徒說聽你們的就是聽我& \multicolumn{1}{p{7.93em}|}{} & \multicolumn{1}{p{5.285em}|}{} &       & 10:16&  \\\midrule


勞苦擔重擔者得安息享安息& \multicolumn{1}{p{7.93em}|}{去耶京途中？} & \multicolumn{1}{p{5.285em}|}{11:28-30} &       &  &  \\ \midrule

對律法師論好鄰舍撒瑪利亞人&       &       &       & 10:25-37 &  \\ \midrule
\end{tabular}%
\end{table}%




\begin{table}[H]
  \centering 
    \begin{tabular}{p{13.07em}|l|l|l|p{5.285em}|r}
    \toprule
    \textbf{事 件} & \multicolumn{1}{p{7.93em}|}{\textbf{地 點}} & \multicolumn{1}{p{5.285em}|}{\textbf{馬太福音}} & \multicolumn{1}{p{5.57em}|}{\textbf{馬可福音}} & \textbf{路加福音} & \multicolumn{1}{p{4.855em}}{\textbf{約翰福音}} \\\midrule 



 \textbf{修殿節對猶太人論永生和父} & \multicolumn{1}{p{7.93em}|}{修殿節所羅門廊} &       &       & \multicolumn{1}{r|}{} & \multicolumn{1}{p{4.855em}}{10:22-30} \\\midrule

\textbf{父在我裡面我在父裡面} & \multicolumn{1}{p{7.93em}|}{修殿節所羅門廊} &       &       & \multicolumn{1}{r|}{} & \multicolumn{1}{p{4.855em}}{10:31-39} \\\midrule

众人信約翰指著耶稣的話是真& \multicolumn{1}{p{7.93em}|}{約旦河外} & \multicolumn{1}{p{5.285em}|}{} && &\multicolumn{1}{p{4.855em}}{10:40-42} \\\midrule

與馬大論馬利亞得上好福份 & \multicolumn{1}{p{7.93em}|}{伯大尼} &       &       & 10:38-42 &  \\ \midrule

耶穌教導門徒禱告得聖靈 &       &       &       & 11:1-13 &  \\\midrule

%趕啞瞎鬼眾人驚其为大衛子孫&\multicolumn{1}{l|}{加利利？}& 12:22-23&& & \\ \midrule
%趕出叫人啞巴的鬼&\multicolumn{1}{l|}{加利利？}& && 11:14& \\ \midrule
%耶穌忙碌，親屬說他癲狂&\multicolumn{1}{l|}{加利利？}&& 3:20-21& & \\ \midrule

被法利賽人文士诬靠鬼王趕鬼&\multicolumn{1}{l|}{加利利？}& 12:24& 3:22& 11:15 & \\ \midrule

文士法利賽人求天上神蹟&\multicolumn{1}{l|}{加利利}&12:38&  & 11:16 & \\ \midrule


撒旦自相攻打必要滅亡&\multicolumn{1}{l|}{}& 12:25-26& 3:23-26& 11:17-18& \\ \midrule

主靠神的靈趕鬼就是神國臨到&\multicolumn{1}{l|}{}& 12:27-28&& 11:19-20& \\ \midrule

先捆住壯士才可搶奪他的家&\multicolumn{1}{l|}{}& 12:29& 3:27& 11:21-22& \\ \midrule

不與主相合就是敵主的&\multicolumn{1}{l|}{}& 12:30& & 11:23& \\ \midrule


七個惡鬼入住打掃乾淨的空屋&\multicolumn{1}{l|}{加利利？}&12:43-45 & & 11:24-26& \\ \midrule


聽道而守之比懷胎乳養者有福&\multicolumn{1}{l|}{加利利?}& &  & 11:27-28& \\ \midrule
約拿的神蹟所羅門的智慧&\multicolumn{1}{l|}{加利利?}& 12:39-42 &  & 11:29-32 & \\ \midrule
論點燈和心裡的光&\multicolumn{1}{l|}{加利利?}& 12:22-45 &  & 11:33-36 & \\ \midrule

法利賽人詫異耶穌飯前不洗手&法利賽人家 &       &       & 11:37-41&  \\ \midrule
論法利賽人律法師的六禍 &法利賽人家&       &       & 11:42-54 &  \\ \midrule

%耶穌與衆人講論? &       &       &       & 12:1-59 &  \\ \midrule
教門徒不要假冒為善,要怕神&       &       &       &12:1-7&  \\ \midrule
褻瀆聖靈不得赦免&       &    12:31-37 &3:28-30      &12:8-12&  \\ \midrule

看果子可知樹惡人心裡發出惡&       &    12:33-35 & &&  \\ \midrule
審判的日子憑話定人為義&       &    12:36-37 & &&  \\ \midrule

教分開家業的人和眾人戒貪心&       &       &       &12:13-15&  \\ \midrule
對眾人論無知財主的自大&       &       &       &12:16-21&  \\ \midrule
教導門徒勿慮衣食&       &       &       &12:22-32&  \\ \midrule
存財寶在天上&       &       &       &12:33-34&  \\ \midrule
忠心警醒等候的僕人&       &       &       &12:35-40&  \\ \midrule
對彼得論忠僕有福&       &       &       &12:41-48&  \\ \midrule

把火丟在地上叫家人紛爭&       &       &       &12:49-53&  \\ \midrule
對眾人說要知道辨別時機&       &       &       &12:54-59&  \\ \midrule
以彼拉多在祭物中摻血論災難&       &       &       &13:1-5&  \\\midrule
以不結實的無花果樹論悔改 &       &       &       &13:6-9 &  \\\midrule
%安息日醫18年駝背女人& 在會堂&       &       & \multicolumn{1}{p{5.285em}|}{13:10-17} & \multicolumn{1}{l}{} \\\midrule
耶穌再論芥種面酵之喻 & 比利亞？   &       &       & \multicolumn{1}{p{5.285em}|}{13:18-21} & \multicolumn{1}{l}{} \\ 
   \bottomrule
    \end{tabular}%
 % \caption{到第四个逾越节前（1）}   
  \label{tab:addlabel}%
\end{table}%



\begin{table}[H]
  \centering
    \begin{tabular}{p{14em}|p{7.93em}|l|l|l|p{4.855em}}
    \toprule
    \textbf{事 件} & \textbf{地 點} & \multicolumn{1}{p{5.285em}|}{\textbf{馬太福音}} & \multicolumn{1}{p{5.57em}|}{\textbf{馬可福音}} & \multicolumn{1}{p{5.285em}|}{\textbf{路加福音}} & \textbf{約翰福音} \\  \midrule


論進神國要努力進窄門&\textcolor{red}{耶京路過的城鄉}&       &       & \multicolumn{1}{p{5.285em}|}{13:22-30} & \multicolumn{1}{l}{} \\\midrule
對法利賽人論希律殺耶穌之事& 比利亞？   &       &       & \multicolumn{1}{p{5.285em}|}{13:31-33} & \multicolumn{1}{l}{} \\\midrule

論耶路撒冷結局,為其悲哀 & 比利亞？   &       &       & \multicolumn{1}{p{5.285em}|}{13:34-35} & \multicolumn{1}{l}{} \\\midrule


警告揀擇首位的客人& 比利亞？   & & & \multicolumn{1}{p{5.285em}|}{14:7-11}& \multicolumn{1}{l}{} \\ \midrule
教訓筵席的主人& 比利亞？   & & & \multicolumn{1}{p{5.285em}|}{14:12-14}& \multicolumn{1}{l}{} \\ \midrule

以大筵席喻猶太人藐視救恩& 比利亞？&&& \multicolumn{1}{p{5.285em}|}{14:15-24} & \multicolumn{1}{l}{} \\ \midrule

從者計算花費,作門徒的代價& 比利亞？   &       &       & \multicolumn{1}{p{5.285em}|}{14:25-35} & \multicolumn{1}{l}{} \\ \midrule


法利賽人文士論主接待稅吏罪人& 比利亞？   & &       & \multicolumn{1}{p{5.285em}|}{15:1-2} & \multicolumn{1}{l}{} \\ \midrule

失羊的比喻 & 比利亞？   &18:12-14       &       & \multicolumn{1}{p{5.285em}|}{15:3-7} & \multicolumn{1}{l}{} \\ \midrule
失錢的比喻 & 比利亞？   & &       & \multicolumn{1}{p{5.285em}|}{15:8-10} & \multicolumn{1}{l}{} \\ \midrule

浪子的比喻 & 比利亞？   &&       & \multicolumn{1}{p{5.285em}|}{15:11-32} & \multicolumn{1}{l}{} \\ \midrule

耶穌設不義聰管家爲喻 & 比利亞？   &       &       & \multicolumn{1}{p{5.285em}|}{16:1-12} & \multicolumn{1}{l}{} \\\midrule
一個僕人不能侍奉兩個主&   &       &       & \multicolumn{1}{p{5.285em}|}{16:13} & \multicolumn{1}{l}{} \\\midrule
责法利賽人自稱為義&   &       &       & \multicolumn{1}{p{5.285em}|}{16:14-17} & \multicolumn{1}{l}{} \\\midrule



法利賽人試探耶穌何故可休妻&猶太境約旦河外 &19:3&10:2& \multicolumn{1}{p{5.285em}|}{} & \multicolumn{1}{l}{} \\\midrule
主答创世记二人成為一體&猶太境約旦河外 &19:4-6&10:6-9& \multicolumn{1}{p{5.285em}|}{} & \multicolumn{1}{l}{} \\\midrule
摩西為何许可給妻子休書&猶太境約旦河外 &19:7-8&10:3-5& \multicolumn{1}{p{5.285em}|}{} & \multicolumn{1}{l}{} \\\midrule

对门徒論休妻另娶的犯姦淫&約旦河外屋裡 & 19:9 &10:10-12& \multicolumn{1}{p{5.285em}|}{16:18} & \multicolumn{1}{l}{} \\\midrule
主以为天国单身解門徒婚姻之惑&約旦河外屋裡  & \multicolumn{1}{p{5.285em}|}{19:10-12} & \multicolumn{1}{p{5.57em}|}{} &       & \multicolumn{1}{l}{} \\\midrule

耶穌論財主與拉撒路 & 比利亞 ？  &       &       & \multicolumn{1}{p{5.285em}|}{16:19-31} & \multicolumn{1}{l}{} \\ \midrule

絆倒人有禍不如頸拴石沉海& 迦百農 ？  &  \multicolumn{1}{p{5.285em}|}{18:6-7} & \multicolumn{1}{p{5.57em}|}{9:42}  & \multicolumn{1}{p{5.285em}|}{17:1-2} & \multicolumn{1}{l}{} \\\midrule

一天七次饒恕懊悔的弟兄& 比利亞 ？  &       &       & \multicolumn{1}{p{5.285em}|}{17:3-4} & \multicolumn{1}{l}{} \\\midrule

對使徒論信心與僕人的本分 & 比利亞 ？  &       &       & \multicolumn{1}{p{5.285em}|}{17:5-10} & \multicolumn{1}{l}{} \\\midrule


對法利賽人論神國在人心裡&去耶京的路上&       &       & \multicolumn{1}{p{5.285em}|}{17:20-21} & \multicolumn{1}{l}{} \\ \midrule

對門徒以挪亞羅得論人子再來&去耶京的路上&       &       & \multicolumn{1}{p{5.285em}|}{17:22-32} & \multicolumn{1}{l}{} \\ \midrule

捨生命者得生命&去耶京的路上&       &       & \multicolumn{1}{p{5.285em}|}{17:33-37} & \multicolumn{1}{l}{} \\ \midrule


以切求的寡婦教導禱告要恒切&去耶京的路上&       &       & \multicolumn{1}{p{5.285em}|}{18:1-8} & \multicolumn{1}{l}{} \\ \midrule

法利賽人稅吏禱告,論自義自卑&去耶京的路上&       &       & \multicolumn{1}{p{5.285em}|}{18:9-14} & \multicolumn{1}{l}{} \\\midrule

\end{tabular}%
\end{table}%


\begin{table}[H]
  \centering
    \begin{tabular}{p{14em}|p{7.93em}|l|l|l|p{4.855em}}
    \toprule
    \textbf{事 件} & \textbf{地 點} & \multicolumn{1}{p{5.285em}|}{\textbf{馬太福音}} & \multicolumn{1}{p{5.57em}|}{\textbf{馬可福音}} & \multicolumn{1}{p{5.285em}|}{\textbf{路加福音}} & \textbf{約翰福音} \\  \midrule

富足官求永生憂愁而去&出來行路的時候&\multicolumn{1}{p{5.285em}|}{19:16-22}&\multicolumn{1}{p{5.285em}|}{10:17-22}& \multicolumn{1}{p{5.285em}|}{18:18-23} & \multicolumn{1}{l}{} \\\midrule

教導門徒財主靠財難進天國&去耶京的路上& \multicolumn{1}{p{5.285em}|}{19:23-26} & \multicolumn{1}{p{5.57em}|}{10:23-27} & \multicolumn{1}{p{5.285em}|}{18:24-27} & \multicolumn{1}{l}{} \\\midrule

對彼得論門徒從主審判十二支派&去耶京的路上& \multicolumn{1}{p{5.285em}|}{19:27-28} & \multicolumn{1}{p{5.57em}|}{} & \multicolumn{1}{p{5.285em}|}{} & \multicolumn{1}{l}{} \\\midrule

對彼得論今世得百倍來世得永生&去耶京的路上& \multicolumn{1}{p{5.285em}|}{19:29-30} & \multicolumn{1}{p{5.57em}|}{10:28-31} & \multicolumn{1}{p{5.285em}|}{18:28-30} & \multicolumn{1}{l}{} \\\midrule

葡萄園雇工各得一錢銀子之喻 && \multicolumn{1}{p{5.285em}|}{20:1-16} &       &       & \multicolumn{1}{l}{} \\\midrule

耶穌對門徒\textcolor{red}{第四次}提及受害事 &去耶京的路上& \multicolumn{1}{p{5.285em}|}{20:17-19} & \multicolumn{1}{p{5.57em}|}{10:32-34} & \multicolumn{1}{p{5.285em}|}{18:31-34} & \multicolumn{1}{l}{} \\\midrule


雅各约翰母亲为子争左右位&  & \multicolumn{1}{p{5.285em}|}{20:20-23} & \multicolumn{1}{p{5.57em}|}{10:35-40} &       & \multicolumn{1}{l}{} \\\midrule
十個門徒聽見惱怒弟兄二人&   & \multicolumn{1}{p{5.285em}|}{20:24} & \multicolumn{1}{p{5.57em}|}{10:41} &       & \multicolumn{1}{l}{} \\\midrule
對門徒論誰願為首必做僕人&   & \multicolumn{1}{p{5.285em}|}{20:25-28} & \multicolumn{1}{p{5.57em}|}{10:42-45} &       & \multicolumn{1}{l}{} \\\midrule
對眾人設交銀與十僕喻神國顯现 & 耶利哥   &       &       & \multicolumn{1}{p{5.285em}|}{19:11-27} & \multicolumn{1}{l}{} \\ \midrule

\end{tabular}%
\end{table}%




\iffalse 
 \begin{table}[htbp]
  \centering
    \begin{tabular}{|p{3.215em}|p{12.215em}|l|r|p{5.145em}|}
    \toprule
    \multicolumn{2}{|p{4.645em}|}{教導內容} & \multicolumn{1}{p{6.645em}|}{馬太福音} & \multicolumn{1}{p{5.355em}|}{馬可福音} & 路加福音 \\
    \midrule
    \multicolumn{2}{|p{13.645em}|}{\textbf{耶穌設立十二門徒}} &       & \multicolumn{1}{p{5.355em}|}{\textbf{3:13-19}} & \textbf{6:12-16} \\
    \midrule
    \multicolumn{1}{|p{3.645em}|}{\multirow{7}{*}{\parbox{0.08\textwidth}{耶穌在山上的教訓}}} & a.論天國子民的性質 & \multicolumn{1}{p{6.645em}|}{5:1-12} &       & \multicolumn{1}{r|}{} \\
\cmidrule{2-5}          & b.論天國子民對世界的影響 & \multicolumn{1}{p{6.645em}|}{5:13-16} &       & \multicolumn{1}{r|}{} \\
\cmidrule{2-5}          & c.論天國子民的義超過律法 & \multicolumn{1}{p{6.645em}|}{5:17-48} &       & \multicolumn{1}{r|}{} \\
\cmidrule{2-5}          & d.論天國子民事奉應有光景 & \multicolumn{1}{p{6.645em}|}{6:1-18} &       & \multicolumn{1}{r|}{} \\
\cmidrule{2-5}          & e.論天國子民對財物的態度 & \multicolumn{1}{p{6.645em}|}{6:19-34} &       & \multicolumn{1}{r|}{} \\
\cmidrule{2-5}          & f.論天國子民對人的原則 & \multicolumn{1}{p{6.645em}|}{7:1-12} &       & \multicolumn{1}{r|}{} \\
\cmidrule{2-5}          & g.論天國子民生活工作根據 & \multicolumn{1}{p{6.645em}|}{7:13-29} &       & \multicolumn{1}{r|}{} \\
 
   \midrule
    \multicolumn{1}{|p{3.645em}|}{\multirow{6}{*}{\parbox{0.08\textwidth}{耶穌在平地教訓衆人}}} & a.論福與禍 &       &       & 6:20-26 \\
\cmidrule{2-5}          & b.論對待仇人 &       &       & 6:27-36 \\
\cmidrule{2-5}          & c.論饒恕人 &       &       & 6:37-38 \\
\cmidrule{2-5}          & d.論樹與果子 &       &       & 6:39-45 \\
\cmidrule{2-5}          & e.論兩等根基 &       &       & 6:46-49 \\
    \bottomrule
    \end{tabular}%
     \caption{到第三個逾越節之前的講論 （1）}
  \label{tab:addlabel}%
\end{table}%
\fi
















\subsubsection{在加利利的講論}
% Table generated by Excel2LaTeX from sheet 'Sheet5'
\begin{table}[htbp]
  \centering
    \begin{tabular}{|p{5.43em}|l|p{6.645em}|p{5.355em}|p{5.145em}|}
    \toprule
    \multicolumn{2}{|p{18.645em}|}{教導內容} & 馬太福音  & 馬可福音  & 路加福音 \\
    \midrule
    \multicolumn{2}{|p{18.645em}|}{\textbf{耶穌論施洗約翰}} & \textbf{11:2-19} & \multicolumn{1}{l|}{} & \textbf{7:18-35} \\
    \midrule
    \multirow{7}{*}{\parbox{0.08\textwidth}{耶穌對法利賽人的講論}} & \multicolumn{1}{p{13.215em}|}{a.耶穌趕出叫人瞎啞的鬼} & 12:22-23 &  & \multicolumn{1}{l|}{11:14} \\
\cmidrule{2-5}    \multicolumn{1}{|l|}{} & \multicolumn{1}{p{13.215em}|}{b.法利賽人妄論耶穌} & \multicolumn{1}{l|}{12:24} & \multicolumn{1}{l|}{3:22} & \multicolumn{1}{l|}{11:15} \\
\cmidrule{2-5}     & \multicolumn{1}{p{13.215em}|}{c.耶穌對他們的反駁} & 12:25-37 & 3:23-30 & 11:17-23 \\
\cmidrule{2-5}    \multicolumn{1}{|l|}{} & d.論這世代的神迹 & 12:38-42 & \multicolumn{1}{l|}{} & \multicolumn{1}{l|}{11:16,19-32} \\
%\cmidrule{3-5}    \multicolumn{1}{|l|}{} &       &  & \multicolumn{1}{l|}{} & 19-32 \\
\cmidrule{2-5}    \multicolumn{1}{|l|}{} & \multicolumn{1}{p{13.215em}|}{e.論這世代的光景} & 12:43-45 & \multicolumn{1}{l|}{} & 11:24-28 \\
\cmidrule{2-5}    \multicolumn{1}{|l|}{} & \multicolumn{1}{p{13.215em}|}{f.論人身上的燈} & \multicolumn{1}{l|}{} & \multicolumn{1}{r|}{} & 11:33-36 \\
    \midrule
    \multicolumn{2}{|p{18.645em}|}{耶穌論親屬} & 12:46-50 & 3:31-35 & 8:19-21 \\
    \bottomrule
    \end{tabular}%
     \caption{到第三個逾越節之前的講論 （2）}
  \label{tab:addlabel}%
\end{table}%

\subsubsection{在迦百農海邊天國的比喻} 
\begin{table}[htbp]
  \centering

    \begin{tabular}{|p{11.645em}|p{5.355em}|l|l|}
    \toprule
     \multicolumn{1}{|p{9.645em}|}{教導內容} & 馬太福音 & 馬可福音 & 路加福音 \\
\midrule
    a.引言  & 13:1-2 & 4:1  & 8:4 \\
    \midrule
    b.撒種的比喻 & 13:3-9 & \multicolumn{1}{p{5.145em}|}{4:2-9} & \multicolumn{1}{p{5.285em}|}{8:5-8} \\
    \midrule
    c.說明用比喻的原因 & 13:10-18 & \multicolumn{1}{p{5.145em}|}{4:10-12} & \multicolumn{1}{p{5.285em}|}{8:9-10} \\
    \midrule
    d.解明撒種比喻 & 13:19-23 & \multicolumn{1}{p{5.145em}|}{4:13-20} & \multicolumn{1}{p{5.285em}|}{8:11-15} \\
    \midrule
    e.論燈與隱瞞的事 & \multicolumn{1}{r|}{} & \multicolumn{1}{p{5.145em}|}{4:21-23} & \multicolumn{1}{p{5.285em}|}{8:16-17} \\
    \midrule
    f.論當留心聽 & \multicolumn{1}{r|}{} & \multicolumn{1}{p{5.145em}|}{4:24-25} & 8:18 \\
    \midrule
    g.種子生長的奧秘 & \multicolumn{1}{r|}{} & \multicolumn{1}{p{5.145em}|}{4:26-29} &  \\
    \midrule
    h.粺子的比喻 & 13:24-30 &       &  \\
    \midrule
    i.芥菜種的比喻 & 13:31-32 & \multicolumn{1}{p{5.145em}|}{4:30-32} &  \\
    \midrule
    j.面酵的比喻 & 13:33-35 & \multicolumn{1}{p{5.145em}|}{4:33-34上} &  \\
    \midrule
  %  k.耶穌進入房子 & 13:36上 &       &  \\
   % \midrule
    k.對門徒解粺子的比喻 & 13:36下-43 & \multicolumn{1}{p{5.145em}|}{4:34下} &  \\
    \midrule
    l.藏寶的比喻 & \multicolumn{1}{l|}{13:44} &       &  \\
    \midrule
    m.尋珠的比喻 & 13:45-46 &       &  \\
    \midrule
    n.撒網的比喻 & 13:47-50 &       &  \\
    \midrule
    o.結論  & 13:51-53 &       &  \\
    \bottomrule
    \end{tabular}%
     \caption{到第三個逾越節之前的講論 （3）}
  \label{tab:addlabel}%
\end{table}%

\subsubsection{其他的講論} 
\begin{table}[htbp]
  \centering
    \begin{tabular}{|p{4.43em}|l|p{5.1em}|l|l|l|p{4.215em}|}
    \toprule
    \multicolumn{2}{|p{10.645em}|}{教導內容} & 地 點   & 馬太福音  & 馬可福音  & 路加福音  & 約翰福音 \\
\midrule
    \multirow{4}{*}{\parbox{0.08\textwidth}{耶穌差遣十二使徒傳道}} & \multicolumn{1}{p{10.215em}|}{a.差遣的動機} & \multirow{4}{*}{加利利} & \multicolumn{1}{p{5.355em}|}{9:36-38，10:1} & 6:7  & \multicolumn{1}{p{4.8em}|}{9:1-2} &  \\
\cmidrule{2-2}\cmidrule{4-7}    \multicolumn{1}{|l|}{} & \multicolumn{1}{p{10.215em}|}{b.十二使徒的名字} & \multicolumn{1}{c|}{} & \multicolumn{1}{p{5.355em}|}{10:2-4} &       &       &   \\
\cmidrule{2-2}\cmidrule{4-7}    \multicolumn{1}{|l|}{} & \multicolumn{1}{p{10.215em}|}{c.打發前的教導和勉勵} & \multicolumn{1}{c|}{} & \multicolumn{1}{p{5.355em}|}{10:5-42} & \multicolumn{1}{p{5.145em}|}{6:8-11} & \multicolumn{1}{p{5.285em}|}{9:3-5} &   \\
\cmidrule{2-2}\cmidrule{4-7}    \multicolumn{1}{|l|}{} & \multicolumn{1}{p{10.215em}|}{d.使徒出去工作} & \multicolumn{1}{c|}{} & 11:1 & \multicolumn{1}{p{5.145em}|}{6:12-13} & 9:6  &   \\
    \midrule
    \multicolumn{2}{|p{18.645em}|}{\textbf{複記施洗約翰被殺經過}} & \textbf{馬格路斯堡} & \multicolumn{1}{p{5.355em}|}{\textbf{14:6-12}} & \multicolumn{1}{p{5.145em}|}{\textbf{6:21-29}} &       & \multicolumn{1}{r|}{} \\
    \midrule
    \multicolumn{2}{|p{18.645em}|}{耶穌論生命的糧} & 迦百農   &       &       &       & 6:22-7:1 \\
    \bottomrule
    \end{tabular}%
     \caption{到第三個逾越節之前的講論 （4）}
  \label{tab:addlabel}%
\end{table}%

\subsection{到第四個逾越節之前的講論}
\subsubsection{住棚節之前的講論}
% Table generated by Excel2LaTeX from sheet 'Sheet2'
\begin{table}[htbp]
  \centering
    \begin{tabular}{|p{5.855em}l|l|l|l|l|p{4.07em}|}
    \toprule
    \multicolumn{2}{|p{12.71em}|}{事 件} & \multicolumn{1}{p{7.93em}|}{地 點} & \multicolumn{1}{p{4.0em}|}{馬太福音} & \multicolumn{1}{p{4.0em}|}{馬可福音} &  \multicolumn{1}{p{4.em}|}{路加福音} & 約翰福音 \\
    \midrule
    \multicolumn{2}{|p{12.71em}|}{耶穌斥責求神迹者} & 近馬加丹 & 16:1-4 & 8:11-13 &       &  \\
    \midrule
    \multicolumn{2}{|p{12.71em}|}{耶穌論潔淨之禮} & \multicolumn{1}{p{4.93em}|}{迦百農} & 15:1-20 & 7:1-23 &       &  \\
    \midrule
    \multicolumn{2}{|p{12.71em}|}{耶穌警告門徒防酵} & \multicolumn{1}{p{4.93em}|}{加利利海} & 16:5-12 & 8:14-21 &       &  \\
    \midrule
    \multicolumn{2}{|p{12.71em}|}{\textbf{耶穌首次提說遭害的事}} & \multicolumn{1}{p{7.93em}|}{該撒利亞腓立比} & 16:21-28 &8:31-9:1 & 9:22-27 & \\
    \midrule
    \multicolumn{2}{|p{12.71em}|}{耶穌下山時談論以利亞} & \multicolumn{1}{p{7.93em}|}{腓立比黑門山} & 17:9-13 & 9:9-13 &       &  \\
    \midrule
    \multicolumn{2}{|p{12.71em}|}{\textbf{耶穌再論受害的事}} & \multicolumn{1}{p{4.93em}|}{加利利} & 17:22-23 &9:30-32 &9:43下-45 &  \\
    \midrule
    \multicolumn{2}{|p{12.71em}|}{耶穌教導門徒謙卑饒恕} & \multicolumn{1}{p{4.93em}|}{迦百農} & 18:1-35 & 9:33-50 & 9:46-50 &  \\
    \midrule
    \multicolumn{1}{|l|}{\multirow{2}{*}{\parbox{0.098\textwidth}{\textbf{耶穌在住棚節時的講論}}}} & \multicolumn{1}{p{7.855em}|}{住棚節時的講論} & \multicolumn{1}{l|}{\multirow{3}{*}{耶路撒冷}} &       &       &       & 7:2-52 \\
\cmidrule{2-2}\cmidrule{4-7}    \multicolumn{1}{|c|}{} & \multicolumn{1}{p{7.855em}|}{耶穌論光與自由} &       &       &       &       & 8:12-59 \\
\cmidrule{1-2}\cmidrule{4-7}    \multicolumn{2}{|p{10.71em}|}{耶穌論好牧人（修殿節）} &       &       &       &       & 10:1-21 \\
    \bottomrule
    \end{tabular}%
     \caption{到第四個逾越節之前的講論（1）}
  \label{tab:addlabel}%
\end{table}%


\subsubsection{到第四個逾越節之前的講論（2）}
\begin{table}[htbp]
  \centering
    \begin{tabular}{|p{6.215em}|p{12.215em}|l|l|p{4.215em}|}
    \toprule
     \multicolumn{2}{|p{18.1em}|}{事 件} & \multicolumn{1}{p{2.3em}|}{地 點} & \multicolumn{1}{p{4.0em}|}{馬太福音} &  \multicolumn{1}{p{4.em}|}{路加福音}  \\
    \midrule
    \multicolumn{2}{|p{10.43em}|}{耶穌論作門徒的代價} & \multicolumn{1}{p{2.215em}|}{途中} & \multicolumn{1}{p{4.215em}|}{8:18-22} & 9:57-62 \\
    \midrule
    \multicolumn{2}{|p{8.43em}|}{耶穌差遣七十人} & \multicolumn{1}{p{2.215em}|}{途中} & \multicolumn{1}{p{4.215em}|}{11:20-24} & 10:1-20 \\
    \midrule
    \multicolumn{2}{|p{8.43em}|}{耶穌歡樂頌神} & \multicolumn{1}{p{2.215em}|}{途中} & \multicolumn{1}{p{4.215em}|}{11:25-30} & 10:21-24 \\
    \midrule
    \multicolumn{2}{|p{8.43em}|}{耶穌論好鄰舍} &       &       & 10:25-37 \\
    \midrule
    \multicolumn{2}{|p{8.43em}|}{耶穌教導門徒禱告} &       &       & 11:1-13 \\
    \midrule
    \multicolumn{2}{|p{8.43em}|}{論法利賽人的禍} &       &       & 11:37-54 \\
    \midrule
    \multicolumn{1}{|c|}{\multirow{7}{*}{\parbox{0.09\textwidth}{耶穌與衆人講論}}} & a.勸門徒提防法利賽人的酵 &       &       & \multicolumn{1}{l|}{12:1} \\
\cmidrule{2-5}          & b.勸門徒不要怕人 &       &       & 12:2-12 \\
\cmidrule{2-5}          & c.對衆人論生命勝過財富 &       &       & 12:13-21 \\
\cmidrule{2-5}          & d.對門徒論神的看顧 &       &       & 12:22-34 \\
\cmidrule{2-5}          & e.對門徒論僕人當儆醒 &       &       & 12:35-48 \\
\cmidrule{2-5}          & f.論主來後地上的光景 &       &       & 12:49-53 \\
\cmidrule{2-5}          & g.論要分辨這世代 &       &       & 12:54-59 \\
    \midrule
    \multicolumn{1}{|c|}{\multirow{2}{*}{\parbox{0.09\textwidth}{耶穌論災難與悔改}}} & a.論悔改與滅亡 &       &       & 13:1-5 \\
\cmidrule{2-5}          & b.論白占土地的樹 &       &       & 13:6-9 \\
    \bottomrule
    \end{tabular}%
      \caption{到第四個逾越節之前的講論（2）}
  \label{tab:addlabel}%
\end{table}%

\subsubsection{到第四個逾越節之前的講論（3）}
\begin{table}[htbp]
  \centering
    \begin{tabular}{|p{5.93em}l|p{4.215em}|r|r|p{5.855em}|}
    \toprule
    \multicolumn{2}{|c|}{\textbf{事 件}} & \textbf{地 點} & \multicolumn{1}{p{5.215em}|}{\textbf{馬太福音}} & \multicolumn{1}{p{5.07em}|}{\textbf{馬可福音}} & \textbf{路加福音} \\
    \midrule
    \multicolumn{2}{|p{12.145em}|}{耶穌再論芥種面酵之喻} & 比利亞   &       &       & 13:18-21 \\
    \midrule
    \multicolumn{2}{|p{12.145em}|}{耶穌論進神國者的資格} & 比利亞   &       &       & 13:22-30 \\
    \midrule
    \multicolumn{2}{|p{12.145em}|}{耶穌論耶路撒冷的結局} & 比利亞   &       &       & 13:31-35 \\
  %  \midrule
  %  \multicolumn{2}{|p{12.145em}|}{耶穌醫治患水臌病的人} & 比利亞   &       &       & 14:1-6 \\
    \midrule
    \multicolumn{2}{|p{12.145em}|}{耶穌設筵席的比喻} & 比利亞   &       &       & 14:7-24 \\
    \midrule
    \multicolumn{2}{|p{12.145em}|}{a.論作客者應注意的事} & \multicolumn{1}{r|}{} &       &       & 14:7-11 \\
    \midrule
    \multicolumn{2}{|p{12.145em}|}{b.論請客者應注意的事} & \multicolumn{1}{r|}{} &       &       & 14:12-14 \\
    \midrule
    \multicolumn{2}{|p{12.145em}|}{c.論神的筵席} & \multicolumn{1}{r|}{} &       &       & 14:15-24 \\
    \midrule
    \multicolumn{2}{|p{12.145em}|}{耶穌要跟從者計算花費} & 比利亞   &       &       & 14:25-35 \\
    \midrule
    \multicolumn{1}{|c|}{\multirow{3}[6]{*}{耶穌論失喪的三個比喻}} & \multicolumn{1}{p{6.215em}|}{a.失落的羊} & \multirow{3}[6]{*}{比利亞} &       &       & 15:1-7 \\
\cmidrule{2-2}\cmidrule{4-6}    \multicolumn{1}{|c|}{} & \multicolumn{1}{p{6.215em}|}{b.失落的錢} & \multicolumn{1}{c|}{} &       &       & 15:8-10 \\
\cmidrule{2-2}\cmidrule{4-6}    \multicolumn{1}{|c|}{} & \multicolumn{1}{p{6.215em}|}{c.失落的人} & \multicolumn{1}{c|}{} &       &       & 15:11-32 \\
    \midrule
    \multicolumn{2}{|p{12.145em}|}{耶穌設不忠的管家爲喻} & 比利亞   &       &       & 16:1-13 \\
    \midrule
    \multicolumn{2}{|p{12.145em}|}{耶穌論財主與拉撒路} & 比利亞   &       &       & 16:14-31 \\
    \midrule
    \multicolumn{2}{|p{12.145em}|}{耶穌論信心與僕人態度} & 比利亞   &       &       & 17:1-10 \\
    \midrule
    \multicolumn{2}{|p{12.145em}|}{耶穌論人子再來的日子} & 以法蓮   &       &       & 17:20-37 \\
    \midrule
    \multicolumn{2}{|p{12.145em}|}{耶穌論禱告要恒切} & 以法蓮   &       &       & 18:1-8 \\
    \midrule
    \multicolumn{2}{|p{12.145em}|}{耶穌論自義的人} & 以法蓮   &       &       & 18:9-14 \\
    \midrule
    \multicolumn{2}{|p{12.145em}|}{耶穌論休妻} & 以法蓮   & \multicolumn{1}{p{5.215em}|}{19:3-12} & \multicolumn{1}{p{5.07em}|}{10:2-12} & \multicolumn{1}{r|}{} \\
    \midrule
    \multicolumn{2}{|p{12.145em}|}{耶穌論財主進神的國} & 以法蓮   & \multicolumn{1}{p{5.215em}|}{19:16-30} & \multicolumn{1}{p{5.07em}|}{10:17-31} & 18:18-30 \\
    \midrule
    \multicolumn{2}{|p{12.145em}|}{耶穌論葡萄園雇工之喻} & 以法蓮   & \multicolumn{1}{p{5.215em}|}{20:1-16} &       & \multicolumn{1}{r|}{} \\
    \midrule
    \multicolumn{2}{|p{12.145em}|}{\textbf{耶穌第三次提及受害事}} & 以法蓮   & \multicolumn{1}{p{5.215em}|}{20:17-19} & \multicolumn{1}{p{5.07em}|}{10:32-34} & 18:31-34 \\
    \midrule
    \multicolumn{2}{|p{12.145em}|}{門徒爭論誰爲大} & 以法蓮   & \multicolumn{1}{p{5.215em}|}{20:20-28} & \multicolumn{1}{p{5.07em}|}{10:35-45} & \multicolumn{1}{r|}{} \\
    \midrule
    \multicolumn{2}{|p{12.145em}|}{耶穌設交銀與僕人之喻} & 耶利哥   &       &       & 19:11-28 \\
    \bottomrule
    \end{tabular}%
     \caption{到第四個逾越節之前的講論（3）}
  \label{tab:addlabel}%
\end{table}%

\subsection{第四個逾越節之間的講論}
\subsubsection{論信心的功效，去耶京的路上（週二）}
\subsubsection{在聖殿的講論（週二）}
\subsubsection{將來之事的講論（週二）}
\subsubsection{逾越節晚餐上的講論（週四）}

\begin{table}[htbp]
  \centering 
    \begin{tabular}{|p{4.57em}l|p{4.43em}|p{4.215em}|p{4.355em}|r|r|}
    \toprule
    \multicolumn{2}{|p{15.285em}|}{\textbf{事 件}} & \textbf{地 點} & \textbf{馬太福音} & \textbf{馬可福音} & \multicolumn{1}{p{4.em}|}{\textbf{路加福音}} & \multicolumn{1}{p{4.em}|}{\textbf{約翰福音}} \\
    \midrule
    \multicolumn{2}{|p{17.285em}|}{耶穌論信心的功效} & 去耶京   & 21:20-22 & 11:20-26 &       &  \\
    \midrule
    \multicolumn{1}{|l|}{\multirow{4}{*}{\parbox{0.08\textwidth}{耶穌論祂的權柄}}} & \multicolumn{1}{p{12.715em}|}{a.對答祭司長的問題} & 耶路撒冷  & 21:23-27 & 11:27-33 & \multicolumn{1}{p{5.145em}|}{20:1-8} &  \\
\cmidrule{2-7}    \multicolumn{1}{|l|}{} & \multicolumn{1}{p{12.715em}|}{b.講兩個兒子聽話的比喻} & \multicolumn{1}{r|}{} & 21:28-32 & \multicolumn{1}{r|}{} &       &  \\
\cmidrule{2-7}    \multicolumn{1}{|l|}{} & \multicolumn{1}{p{12.715em}|}{c.講惡園戶的比喻} & \multicolumn{1}{r|}{} & 21:33-46 & 12:1-12 & \multicolumn{1}{p{5.145em}|}{20:9-19} &  \\
\cmidrule{2-7}    \multicolumn{1}{|l|}{} & \multicolumn{1}{p{12.715em}|}{d.講設筵席的比喻} & \multicolumn{1}{r|}{} & 22:1-14 & \multicolumn{1}{r|}{} &       &  \\
    \midrule
    \multicolumn{2}{|p{17.285em}|}{耶穌論納稅的事} & 耶路撒冷  & 22:15-22 & 12:13-17 & \multicolumn{1}{p{5.145em}|}{20:20-26} &  \\
    \midrule
    \multicolumn{2}{|p{17.285em}|}{耶穌論復活的事} & 耶路撒冷  & 22:23-33 & 12:18-27 & \multicolumn{1}{p{5.145em}|}{20:27-38} &  \\
    \midrule
    \multicolumn{2}{|p{17.285em}|}{耶穌論最大的誡命} & 耶路撒冷  & 22:34-40 & 12:28-34 & \multicolumn{1}{p{5.145em}|}{20:39-40} &  \\
    \midrule
    \multicolumn{2}{|p{17.285em}|}{耶穌論基督與大衛關係} & 耶路撒冷  & 22:41-46 & 12:35-37 & \multicolumn{1}{p{5.145em}|}{20:41-44} &  \\
    \midrule
    \multicolumn{2}{|p{17.285em}|}{耶穌論假冒爲善者光景} & 耶路撒冷  & 23:1-39 & 12:38-40 & \multicolumn{1}{p{5.145em}|}{20:45-47} &  \\
    \midrule
    \multicolumn{2}{|p{17.285em}|}{耶穌稱讚寡婦的奉獻} & 耶路撒冷  & \multicolumn{1}{r|}{} & 12:41-44 & \multicolumn{1}{p{5.145em}|}{21:1-4} &  \\
    \midrule
    \multicolumn{2}{|p{17.285em}|}{耶穌最後公開的講論} & 耶路撒冷  & \multicolumn{1}{r|}{} & \multicolumn{1}{r|}{} &       & \multicolumn{1}{p{5.93em}|}{12:37-50} \\
    \midrule
    \multicolumn{1}{|l|}{\multirow{6}{*}{\parbox{0.08\textwidth}{耶穌論要來之事}}}  & \multicolumn{1}{p{12.715em}|}{a.預言聖殿被拆} & \multirow{6}[12]{*}{橄欖山} & 24:1-14 & 13:1-13 & \multicolumn{1}{p{5.145em}|}{21:5-19} &  \\
\cmidrule{2-2}\cmidrule{4-7}    \multicolumn{1}{|l|}{} & \multicolumn{1}{p{12.715em}|}{b.再臨的預兆} & \multicolumn{1}{c|}{} & 24:15-31 & 13:14-27 & \multicolumn{1}{p{5.145em}|}{21:20-28} &  \\
\cmidrule{2-2}\cmidrule{4-7}    \multicolumn{1}{|l|}{} & \multicolumn{1}{p{12.715em}|}{d.儆醒的勸勉} & \multicolumn{1}{c|}{} & 24:32-51 & 13:28-37 & \multicolumn{1}{p{5.145em}|}{21:29-36} &  \\
\cmidrule{2-2}\cmidrule{4-7}    \multicolumn{1}{|l|}{} & \multicolumn{1}{p{12.715em}|}{d.十個童女的比喻} & \multicolumn{1}{c|}{} & 25:1-13 & \multicolumn{1}{r|}{} &       &  \\
\cmidrule{2-2}\cmidrule{4-7}    \multicolumn{1}{|l|}{} & \multicolumn{1}{p{12.715em}|}{e.分銀子與算帳的比喻} & \multicolumn{1}{c|}{} & 25:14-30 & \multicolumn{1}{r|}{} &       &  \\
\cmidrule{2-2}\cmidrule{4-7}    \multicolumn{1}{|l|}{} & \multicolumn{1}{p{12.715em}|}{f.分別綿羊與山羊的比喻} & \multicolumn{1}{c|}{} & 25:31-46 & \multicolumn{1}{r|}{} &       &  \\
    \midrule
    \multicolumn{2}{|p{17.285em}|}{\textbf{耶穌第四次提及受害事}} & 耶路撒冷  & 26:1-2 & \multicolumn{1}{r|}{} &       &  \\
    \midrule
    \multicolumn{2}{|p{17.285em}|}{耶穌爲門徒洗腳} & 耶路撒冷  & \multicolumn{1}{r|}{} & \multicolumn{1}{r|}{} &       & \multicolumn{1}{p{5.93em}|}{13:1-20} \\
    \midrule
    \multicolumn{2}{|p{17.285em}|}{耶穌賜新命令} & 耶路撒冷  & \multicolumn{1}{r|}{} & \multicolumn{1}{r|}{} &       & \multicolumn{1}{p{5.93em}|}{13:34-35} \\
    \midrule
    \multicolumn{1}{|l|}{\multirow{4}{*}{\parbox{0.08\textwidth}{耶穌最後的談話}}} & \multicolumn{1}{p{12.715em}|}{a.論祂去的目的} & \multirow{4}[8]{*}{耶路撒冷} & \multicolumn{1}{r|}{} & \multicolumn{1}{r|}{} &       & \multicolumn{1}{p{5.93em}|}{14:1-31} \\
\cmidrule{2-2}\cmidrule{4-7}    \multicolumn{1}{|l|}{} & \multicolumn{1}{p{12.715em}|}{b.葡萄樹與門徒彼此相愛} & \multicolumn{1}{c|}{} & \multicolumn{1}{r|}{} & \multicolumn{1}{r|}{} &       & \multicolumn{1}{p{5.93em}|}{15:1-27} \\
\cmidrule{2-2}\cmidrule{4-7}    \multicolumn{1}{|l|}{} & \multicolumn{1}{p{12.715em}|}{c.論聖靈的降臨} & \multicolumn{1}{c|}{} & \multicolumn{1}{r|}{} & \multicolumn{1}{r|}{} &       & \multicolumn{1}{p{5.93em}|}{16:1-16} \\
\cmidrule{2-2}\cmidrule{4-7}    \multicolumn{1}{|l|}{} & \multicolumn{1}{p{12.715em}|}{d.論奉主名禱告} & \multicolumn{1}{c|}{} & \multicolumn{1}{r|}{} & \multicolumn{1}{r|}{} &       & \multicolumn{1}{p{5.93em}|}{16:17-33} \\
    \midrule
    \multicolumn{2}{|p{17.285em}|}{耶穌爲門徒代禱} & 耶路撒冷  & \multicolumn{1}{r|}{} & \multicolumn{1}{r|}{} &       & \multicolumn{1}{p{5.93em}|}{17:1-26} \\
    \midrule
    \multicolumn{2}{|p{17.285em}|}{耶穌最後的命令} & 加利利山  & 28:18-20 & 16:15-18 &       &  \\
    \bottomrule
    \end{tabular}%
     \caption{第四個逾越節期間的講論}
  \label{tab:addlabel}%
\end{table}%
